\documentclass{article}

\usepackage[centertags,fleqn]{amsmath}
\setlength{\mathindent}{0pt} % Get equations completely aligned on the left.
\usepackage{amssymb,amsfonts,xspace}
\usepackage{booktabs,multirow}
\usepackage{longtable}
\usepackage{syntax}

\usepackage{todonotes}
\newcommand{\todoi}[1]{\todo[inline]{todo: #1}}

\usepackage{tikz}

\usepackage{amsthm}
\usepackage{mathtools}
%\usepackage{thmtools}
%\usepackage{thm-restate}
%% \declaretheorem[name=Theorem,numberwithin=section]{theorem}
%% \declaretheorem[name=Lemma,numberwithin=section]{lemma}
%% \declaretheorem[name=Proposition,numberwithin=section]{proposition}
%% \declaretheorem[name=Corollary,numberwithin=section]{corollary}
%% \declaretheorem[name=Conjecture,numberwithin=section]{conjecture}
\newtheorem{theorem}{Theorem}
\newtheorem{question}{Question}
\newtheorem{observation}{Observation}
%\newtheorem{theorem}{Theorem}
\newtheorem{definition}{Definition}
%\newtheorem{lemma}{Lemma}
\newtheorem{proposition}{Proposition}
%\newtheorem{corollary}{Corollary}
\newtheorem{example}[theorem]{Example}
%\newtheorem{conjecture}{Conjecture}
\newtheorem{fact}{Fact}

\usepackage{acronym}
\acrodef{CSP}{Constraint Satisfaction Problem}

\newcommand{\suchthat}{\mid}
\newcommand{\eqdef}{\coloneqq}
\newcommand{\booleans}{\mathbb{B}}
\newcommand{\tuple}[1]{\left\langle#1\right\rangle}
\newcommand{\bigo}{\mathcal{O}}
\newcommand{\abs}[1]{\left\lvert#1\right\rvert}
\newcommand{\integers}{\mathbb{Z}}
\newcommand{\nonnegintegers}{\mathbb{Z}_{\geq 0}}
\newcommand{\floor}[1]{\left\lfloor#1\right\rfloor}
\newcommand{\ceil}[1]{\left\lceil#1\right\rceil}

\newcommand{\valuations}{\mathit{Val}}
\newcommand{\domainof}[1]{D(#1)}
\newcommand{\determination}[1]{G(#1)}
\newcommand{\scope}[1]{S(#1)}
\newcommand{\outputs}[1]{O(#1)}
\newcommand{\inputs}[1]{I(#1)}
\newcommand{\variables}{X}
\newcommand{\constraints}{\mathcal{C}}

\begin{document}

%\section{Fact format}

\section{Syntax}

\subsection{Basic terms}
\begin{grammar}

<type> ::=
~\alt `bool' | `float' | `int' | `none' | `string' | `symbol'
\alt `function' | `set' | `multimap'

<bool> ::= `true' | `false'

<int> ::= any integer literal

<string> ::= any string literal (enclosed in quotes)

<name> ::= any suitable string

<term> ::= <int> | <string> | <symbol>

<terms> ::= <term> | <term> `,' <terms>

<symbol> ::= <name> | <name> `(' <terms> `)'

<term-list> ::= `(' `)' | `(' <term> `,' <term-list> `)'

%<constant> ::= `bool' `(' <bool> `)' \alt `int' `(' <int> `)' \alt `float' `(' <string> `)' \alt `string `(' <string> `)' \alt `symbol `(' <term> `)'

\end{grammar}

\subsection{Building expression terms}
\begin{grammar}

<val> ::= `bad' | `val' `(' <type> `,' <term> `)'

<operator> ::=
~\alt `python' `(' <string> `)'
\alt <lambda-expr>
\alt <variable>
\alt <bool-operator>
\alt <conditional-operator>
\alt <float-operator>
\alt <int-operator>
\alt <multimap-operator>
\alt <set-operator>
\alt <string-operator>

<eq-operator> ::= `eq' | `neq'

<comp-operator> ::= <eq-operator> | `leq' | `lt' | `geq' | `gt'

<bool-operator> ::=
~\alt <eq-operator> | `conj' | `disj' | `leqv'
\alt `limp' | `lnot' | `lxor' | `snot' | `wnot'

<conditional-operator> ::= `if' | `getOrElse' | `ite' | `hasValue'

<float-operator> ::=
~\alt <comp-operator> | `sqrt' | `cos' | `sin' | `tan'
\alt `acos' | `asin' | `atan' | `abs' | `minus'
\alt `add' | `sub' | `mult' | `int_div' | `float_div' | `pow' | `floor' | `ceil'
\alt `max' | `min'

<int-operator> ::=
~\alt <comp-operator> | `add' | `sub' | `mult'
\alt `int_div' | `float_div' | `pow' | `abs' | `minus' | `max' | `min'

<multimap-operator> ::=
~\alt <eq-operator> | `find' | `multimap_fold' | `multimap_isin' | `multimap_make'
\alt `countKeys' | `countEntries' |`sumIntEntries'
\alt `maxEntries' | `minEntries'

<set-operator> ::=
~\alt <eq-operator> | `union' | `inter' | `diff' | `subset' | `set_make'
\alt `set_isin' | `set_notin' | `cardinality' | `set_fold'

<string-operator> ::= <eq-operator> | `concat' | `length'

<lambda-expr> ::= `lambda' `(' <term-list> `,' <expression> `)'

<tuple-expr> ::= `(' `)' | `(' <expressions> `)'

<expression> ::=
~\alt <val>
\alt `variable' `(' <term> `)'
\alt `operation' `(' <operator> `,' <expression-list> `)'
\alt <lambda-expr>
\alt `python' `(' <string> `)'
\alt <tuple-expr>

<expression-list> ::= `(' `)' | `(' <expression> `,' <expression-list> `)'

<expressions> ::= <expression> | <expression> `,' <expressions>

\end{grammar}

\subsection{Building statement terms}

\begin{grammar}
<statement> ::=
~\alt `assert' `(' <expression> `)'
\alt `assign' `(' <term> `,' <expression> `)'
\alt `if' `(' <expression> `,' <statement> `,' <statement> `)'
\alt `noop'
\alt `statement_python' `(' <string> `)'
\alt `seq2' `(' <statement> `,' <statement> `)'
\alt `while' `(' <int> `,' <expression> `,' <statement> `)'
\end{grammar}

\subsection{Building facts and declarations}

\begin{grammar}
<domain> ::= `definition' | `boolDomain' | `fromFacts' | `open' | `set' | `multimap'

<variable-atom> ::=
~\alt `variable_assign' `(' <term> `,' <expression> `)'
\alt `variable_assign' `(' <term> `,' <term> `,' <expression> `)'
\alt `variable_choice' `(' <term> `,' <expression> `)'
\alt `variable_choice' `(' <term> `,' <term> `,' <expression> `)'
\alt `variable_declare' `(' <term> `,' <domain> `)'
\alt `variable_declare' `(' <term> `,' <term> `,' <domain> `)'
\alt `variable_default' `(' <term> `,' <expression> `)'
\alt ...
\alt `variable_default' `(' <term> `,' <term> `,' <expression> `,' <expression> `,' <int> `)'
\alt `variable_define' `(' <term> `,' <expression> `)'
\alt `variable_define' `(' <term> `,' <term> `,' <expression> `)'
\alt `variable_domain' `(' <term> `,' <expression> `)'

<multimap-atom> ::=
~\alt `multimap_assign' `(' <term> `,' <expression> `,' <expression> `)'
\alt `multimap_assign' `(' <term> `,' <term> `,' <expression> `,' <expression> `)'

<set-atom> ::=
~\alt `set_assign' `(' <term> `,' <expression> `)'
\alt `set_assign' `(' <term> `,' <term> `,' <expression> `)'
\alt `set_baseDomain' `(' <term> `,' <expression> `)'
\alt `set_baseDomain' `(' <term> `,' <term> `,' <expression> `)'

<execution-atom> ::=
~\alt `execution_declare' `(' <term> `,' <statement> `,' <term-list> `,' <term-list> `)'
\alt `execution_declare' `(' <term> `,' <term> `,' <statement> `,' <term-list> `,' <term-list> `)'
\alt `execution_run' `(' <term> `)'
\alt `execution_run' `(' <term> `,' <term> `)'

<optimize-atom> ::=
~\alt `optimize_maximizeSum' `(' <expression> `,' <term> `)'
\alt `optimize_maximizeSum' `(' <expression> `,' <term> `,' <expression> `)'
\alt `optimize_maximizeSum' `(' <term> `,' <expression> `,' <term> `,' <expression> `)'
\alt `optimize_precision' `(' <expression>  `)'
\alt `optimize_precision' `(' <expression> `,' <expression> `)'

<preference-atom> ::=
~\alt `preference_maximizeScore'
\alt `preference_holds' `(' <expression> `)'
\alt `preference_holds' `(' <expression> `,' <int> `)'
\alt `preference_holds' `(' <term> `,' <expression> `,' <int> `)'
\alt `preference_variableValue' `(' <term> `,' <expression> `)'
\alt `preference_variableValue' `(' <term> `,' <expression> `,' <int> `)'
\alt `preference_variableValue' `(' <term> `,' <term> `,' <expression> `,' <int> `)'

<warning-control-atom> ::=
~\alt `warning_forbid' `(' <term> `)'
\alt `warning_forbid' `(' <term> `,' <term> `)'
\alt `warning_ignore' `(' <term> `)'
\alt `warning_ignore' `(' <term> `,' <term> `)'

<atom> ::=
~\alt `ensure' `(' <expression> `)'
\alt `ensure' `(' <term> `,' <expression> `)'
\alt `bool_evaluate' `(' <expression> `)'
\alt `bool_evaluate' `(' <term> `,' <expression> `)'
\alt `evaluate' `(' <expression> `)'
\alt `evaluate' `(' <term> `,' <expression> `)'
\alt <variable-atom>
\alt <multimap-atom>
\alt <set-atom>
\alt <execution-atom>
\alt <optimize-atom>
\alt <preference-atom>
\alt <warning-control-atom>
\end{grammar}

\subsection{Result fact format}
\begin{grammar}

<expression-warning> ::=
~\alt `notImplemented'
\alt `pythonError'
\alt `syntaxError'
\alt `zeroDivisionError'

<preference-warning> ::=
~\alt `unsupported'

<statement-warning> ::=
~\alt `evaluatorError'
\alt `notImplemented'
\alt `pythonError'

<type-warning> ::=
~\alt `failed_operation'

<variable-warning> ::=
~\alt `badValue'
\alt `emptyDomain'
\alt `multipleDeclarations'
\alt `multipleDefinitions'
\alt `undeclared'
\alt `confusingName'

<warning-symbol> ::=
~\alt `expression' `(' <expression-warning> `)'
\alt `otherError'
\alt `preference' `(' <preference-warning> `)'
\alt `propagator'
\alt `statement' `(' <statement-warning> `)'
\alt `type' `(' <type-warning> `)'
\alt `variable' `(' <variable-warning> `)'

<atom> ::=
~\alt `value' `(' <term> `,' <val> `)'
\alt `evaluated' `(' <operator> `,' <expression-list>  `,' <val> `)'
\alt `set_value' `(' <term> `,' <val> `)'
\alt `multimap_value' `(' <term> `,' <val> `,' <val> `)'
\alt `preference_score' `(' <int> `)'
\alt `warning' `(' <warning-symbol> `,' <term-list> `,' <term> `)'
\end{grammar}

\subsection*{Examples}
The following two assignments and two checks
\begin{align*}
\pi &\eqdef 3.14159\\
y &\eqdef (x+2) \times \pi \\
& y / 3 \geq 10
& y / 3 \leq 15
\end{align*}
\begin{verbatim}
variable_define(pi,val(float,float("3.14159"))).
variable_define(y,operation(mult,
                           (operation(add,
                                      (variable(x),
                                      (val(int,2),()))),
                            (variable(pi),())))).
ensure(operation(geq,
                (operation(float_div,
                          (variable(y),
                          (val(int,3),()))),
                (val(int,10),())))).
ensure(operation(leq,
                (operation(float_div,
                          (variable(y),
                          (val(int,3),()))),
                (val(int,15),())))).
\end{verbatim}

\paragraph{Other examples}

\begin{verbatim}
variable_define(py_named_function,operation(python("math.cos"),(val(float,float("0.1")),()))).
variable_define(py_lambda,operation(python("lambda x,y: x+y*y"),
                                   (val(float,float("2.0")),
                                   (val(float,float("0.1")),())))).

ensure(operation(eq,
                (variable(v1),
                (variable(v2),
                ()))))
\end{verbatim}

\paragraph{getOrElse function}

\begin{verbatim}
variable_define(x0,val(int,0)).
variable_define(x1,operation(if,(variable(cond),(val(int,1),())))).
variable_define(y0,operation(getOrElse,(variable(x0),(variable(x1),())))).
\end{verbatim}


\subsection{Notes}

Difference with, say, Python:
\begin{itemize}
\item \texttt{a and False} has value False even if \texttt{a} is none.
\end{itemize}


\end{document}
